% -------------------------------------------------------------------------
% WBS ID: RES-DAT-SCN
% Deliverable: Scenario and Severity Definition
% Status: In progress
% Evidence status: Tohoku corridor scenario baseline established
% Review status: Not reviewed
% Last updated: 2026-07-19
% Dependencies: RES-DAT-OBS, RES-DAT-SRC, RES-DAT-GEO
% Downstream interfaces: RES-DAT-CAS, RES-ML-SUR
% -------------------------------------------------------------------------

\section{RES-DAT-SCN --- Scenario and Severity Definition}
\label{sec:res-dat-scn}

\subsection{Purpose and Scope}
\label{sec:res-dat-scn-purpose}

This Deliverable defines one primary historical tsunami event and controlled
physical variations around that event for impact testing. Unrelated tsunami
events shall not be added to the initial project scope.

The active scenario baseline is the 2011 Tōhoku event with two
regional-to-local trajectories: epicentre to the Kamaishi region and
epicentre to the Sendai coastal plain. Scenario definition owns the
evidence-based corridor footprint; it does not define arbitrary drawing
operations or solver implementation.

\subsection{Research Questions}
\label{sec:res-dat-scn-research-questions}

The scenario decision asks:

\begin{enumerate}
  \item How are the two active regional-to-local corridors defined from
    event and target coordinates?
  \item Which evidence determines corridor width, offshore length,
    inland length and sponge/relaxation margins?
  \item Which scenario changes are active impact-test variations and
    which are deferred geometry or ML-generation tasks?
\end{enumerate}

\subsection{Terminology and Definitions}
\label{sec:res-dat-scn-definitions}

% TODO: Define the technical terminology, symbols, quantities and
% classifications used in this Deliverable.

\subsection{Literature Evidence}
\label{sec:res-dat-scn-literature-evidence}

% TODO: Record evidence from peer-reviewed papers, authoritative datasets,
% standards, technical reports and primary documentation.
%
% Do not create a detached literature-summary list. Explain how each source
% supports, contradicts or limits a project decision.

\subsection{Governing Relations and Quantitative Definitions}
\label{sec:res-dat-scn-governing-relations}

% TODO: Record relevant equations, dimensionless groups, empirical models,
% extraction rules or quantitative definitions.
%
% Use align environments for displayed equations.
% Every displayed equation must have a unique \label{eq:...}.
% Do not place punctuation after displayed equations.

\subsubsection{Corridor Coordinate Basis}
\label{sec:res-dat-scn-corridor-coordinate-basis}

Base geographic coordinates are longitude/latitude. The event point is:

\begin{align}
  P_{\mathrm{e}}
  &=
  \left(
    \lambda_{\mathrm{e}},
    \phi_{\mathrm{e}}
  \right)
  \label{eq:res-dat-scn-event-point}
\end{align}

For a target coast point
$P_{\mathrm{t}}=(\lambda_{\mathrm{t}},\phi_{\mathrm{t}})$, a local
metric displacement vector shall be constructed after projection or
local tangent-plane conversion:

\begin{align}
  \Delta\mathbf{r}
  &=
  \left[
    \Delta E,\,
    \Delta N
  \right]^{\mathsf{T}}
  \label{eq:res-dat-scn-target-displacement}
\end{align}

The signed rotation from the latitude/north axis to the corridor
centreline is:

\begin{align}
  \theta
  &=
  \operatorname{atan2}
  \left(
    \Delta E,\,
    \Delta N
  \right)
  \label{eq:res-dat-scn-corridor-bearing-angle}
\end{align}

where positive angle is east of north under the selected projected
coordinate convention. The rotated local basis is:

\begin{align}
  \hat{\mathbf{e}}_{x'}
  &=
  \left[
    \sin\theta,\,
    \cos\theta
  \right]^{\mathsf{T}}
  \label{eq:res-dat-scn-corridor-centreline-basis}
  \\
  \hat{\mathbf{e}}_{y'}
  &=
  \left[
    \cos\theta,\,
    -\sin\theta
  \right]^{\mathsf{T}}
  \label{eq:res-dat-scn-corridor-width-basis}
\end{align}

For any projected point $\mathbf{r}$ measured from $P_{\mathrm{e}}$,
the corridor coordinates are:

\begin{align}
  x'
  &=
  \Delta\mathbf{r}
  \cdot
  \hat{\mathbf{e}}_{x'}
  \label{eq:res-dat-scn-corridor-x-prime}
  \\
  y'
  &=
  \Delta\mathbf{r}
  \cdot
  \hat{\mathbf{e}}_{y'}
  \label{eq:res-dat-scn-corridor-y-prime}
\end{align}

Here $x'$ is the corridor centreline/propagation direction and $y'$ is
the corridor-width direction. A constant-width baseline corridor is:

\begin{align}
  \mathcal{C}
  &=
  \left\{
    \mathbf{r}:
    -L_{\mathrm{pre}}
    \leq
    x'
    \leq
    L_{\mathrm{inland}},
    \quad
    \left|y'\right|
    \leq
    \frac{W}{2}
  \right\}
  \label{eq:res-dat-scn-constant-width-corridor}
\end{align}

where $W$ is corridor width, $L_{\mathrm{pre}}$ is offshore/upstream
pre-impact length and $L_{\mathrm{inland}}$ is inland/downstream
coverage. Sponge or relaxation margins are added outside
$\mathcal{C}$ according to `RES-DAT-GEO`, `RES-MOD-IBC` and
`RES-NUM-IBC`.

Decision trace: Source role: supports coordinate-transform and
corridor-polygon generation for the data/model interface. Supporting
citations: the event coordinate is anchored to
\textcite{usgs2011tohokuEvent}; corridor width and inland/offshore
extent shall be supported by topobathymetric sources
\parencite{jodcJegg500,gsiDemDownload,gebco2026grid} and severe-impact
observations \parencite{mori2012nationwide,noaaNceiHistoricalTsunami}.
Limitation: the angles below use current coordinate proxies and must be
recomputed from the final GIS-selected target coordinates.

\subsection{Required Data Variables and Units}
\label{sec:res-dat-scn-required-data-variables-units}

Required scenario variables are
$P_{\mathrm{e}}$, $P_{\mathrm{t}}$, $\theta$, $W$,
$L_{\mathrm{pre}}$, $L_{\mathrm{inland}}$, upstream/downstream/lateral
sponge widths, target-region identifier, observation-footprint polygon,
defence-geometry footprint and source-data version.

\subsection{Candidate Data Sources}
\label{sec:res-dat-scn-candidate-data-sources}

% TODO: Complete this RES-DAT Domain-specific subsection.

\subsection{Spatial and Temporal Coverage}
\label{sec:res-dat-scn-spatial-temporal-coverage}

The two active trajectories are:

\begin{description}
  \item[$T_{\mathrm{K}}$] epicentre $\rightarrow$ Kamaishi region,
    using the current coordinate proxy
    $\theta_{\mathrm{K}}\approx -21.1^{\circ}$.
  \item[$T_{\mathrm{S}}$] epicentre $\rightarrow$ Sendai coastal
    plain, using the current coordinate proxy
    $\theta_{\mathrm{S}}\approx -97.4^{\circ}$.
\end{description}

The angles are not final GIS bearings. They shall be recomputed from
the accepted target coordinates, selected projection and corridor
origin before clipping or mesh generation.

\subsection{Coordinate and Datum Requirements}
\label{sec:res-dat-scn-coordinate-datum-requirements}

% TODO: Complete this RES-DAT Domain-specific subsection.

\subsection{Data Processing and Transformation}
\label{sec:res-dat-scn-data-processing-transformation}

% TODO: Complete this RES-DAT Domain-specific subsection.

\subsection{Uncertainty, Provenance and Licensing}
\label{sec:res-dat-scn-uncertainty-provenance-licensing}

% TODO: Complete this RES-DAT Domain-specific subsection.

\subsection{Data Acceptance Criteria}
\label{sec:res-dat-scn-data-acceptance-criteria}

A corridor definition is accepted only when:

\begin{itemize}
  \item final $P_{\mathrm{e}}$ and $P_{\mathrm{t}}$ coordinates are
    recorded with source and CRS metadata;
  \item $\theta$ is recomputed from the final GIS coordinates;
  \item $W$, $L_{\mathrm{pre}}$, $L_{\mathrm{inland}}$ and sponge
    widths are justified by source footprint, severe-impact
    observations, defence geometry, bathymetry/topography coverage and
    numerical relaxation requirements;
  \item observation and terrain datasets fully cover the corridor plus
    buffer;
  \item constant-width corridor adequacy is documented for each target.
\end{itemize}

\subsection{Candidate Approaches}
\label{sec:res-dat-scn-candidate-approaches}

The active scenario approach is a validated event baseline plus bounded local
variations in wave height, inundation depth, velocity, direction, duration and
selected severity scaling. These variations support robustness and
reinforcement screening within the event-conditioned domain.

The baseline corridor is constant-width. A narrowing or bay-following
corridor is deferred unless GIS evidence shows that the narrowing
represents true physical bay, coastline or bathymetric constraints
rather than a drawing convenience.

\subsection{Critical Comparison}
\label{sec:res-dat-scn-critical-comparison}

% TODO: Compare candidates using physically and project-relevant criteria.
% Do not select an approach solely on popularity or implementation ease.

The constant-width corridor is selected for the baseline because it
creates a reproducible clipping and replay-boundary interface. It is
limited where a ria bay, harbour mouth, river channel or coastal plain
controls the flow in a way that a rectangular corridor would
misrepresent. In those cases the physical constraint shall be encoded
from GIS coastline/topography evidence, not manually narrowed for
visual neatness.

\subsection{Assumptions and Applicability}
\label{sec:res-dat-scn-assumptions}

% TODO: State assumptions and the conditions under which they are valid.

\subsection{Limitations and Failure Conditions}
\label{sec:res-dat-scn-limitations}

% TODO: State where the evidence, model or selected approach ceases to be
% reliable or applicable.

\subsection{Project-Specific Conclusions}
\label{sec:res-dat-scn-conclusions}

The initial project is a single-event reconstruction and controlled
impact-testing study, not a multi-event generalisation study.

The Tōhoku/local comparison is explicitly two-corridor: Kamaishi for a
Sanriku/ria and defence-interaction case, and Sendai coastal plain for
flat-plain inundation comparison. This pairing supports barrier-class
comparison within one event while avoiding claims of cross-event
generalisability.

\subsection{Selected Approach and Rationale}
\label{sec:res-dat-scn-selected-approach}

Use two evidence-defined, constant-width baseline corridors:
$T_{\mathrm{K}}$ and $T_{\mathrm{S}}$. Corridor parameters are data
products owned by RES-DAT, not solver constants. Source role: supports
regional-to-local data extraction and replay-boundary generation.
Downstream handoff: Software shall generate corridor polygons,
clip raster/vector data, interpolate terrain, extract observation
subsets and write replay-boundary definitions from the same manifest.

\subsection{Unresolved Decisions}
\label{sec:res-dat-scn-unresolved-decisions}

Open decisions are final GIS-selected target coordinates, final
projection, corridor widths, upstream/downstream lengths, sponge
widths, observation-footprint inclusion rules and whether a
physically-constrained non-rectangular local boundary is required near
Kamaishi Bay.

\subsection{Requirements and Handoffs}
\label{sec:res-dat-scn-requirements}

Software shall start the corridor workflow in parallel with remaining
research: coordinate-transform module, corridor-polygon generation,
manifest-driven clipping, local inlet section extraction and
replay-boundary metadata. `RES-VER-TST` shall use the same corridor
parameters in interface-location and one-way-versus-coupled tests.

\subsection{Deliverable Completion Record}
\label{sec:res-dat-scn-completion-record}

% TODO: Record whether the Deliverable is:
%
% - Not started
% - In progress
% - Draft complete
% - Reviewed
% - Accepted
%
% Acceptance requires:
%
% - the research questions to be answered;
% - claims to be supported by appropriate evidence;
% - assumptions and limitations to be explicit;
% - the project-specific conclusion to be stated;
% - downstream requirements to be traceable;
% - unresolved decisions to be closed or formally deferred.
